\clearpage
\section{ПРОЕКТИРОВАНИЕ МЕТОДА И ЭКСПЕРИМЕНТАЛЬНОГО СТЕНДА}

\subsection{Общая архитектура исследования}

В данной работе реализуется не абстрактная схема сравнения методов, а конкретное программное решение, предназначенное для прогнозирования и оптимизации торговых решений на исторических биржевых котировках. Архитектура решения строится как последовательный pipeline обработки данных и принятия решений, в котором каждый модуль имеет четко определенные входы и выходы. Такой подход необходим для того, чтобы все сравниваемые методы находились в одинаковых условиях: использовали одну и ту же исходную информацию, одинаковую схему разбиения данных и единый набор метрик.

Архитектурно программное решение состоит из следующих взаимосвязанных компонентов:
\begin{enumerate}[leftmargin=1.25cm,label=\arabic*)]
\item модуль загрузки и предобработки биржевых данных, формирующий исходный набор OHLCV-наблюдений;
\item модуль формирования признаков, преобразующий ряд цен в набор технических и статистических характеристик;
\item модуль baseline-стратегий и ML-моделей, генерирующий торговые сигналы на основе фиксированных правил или прогноза направления движения;
\item среда обучения с подкреплением, описывающая состояние рынка, действия агента и правила обновления портфеля;
\item RL-агент, обучаемый принимать решения \textit{buy}, \textit{sell}, \textit{hold} на основе наблюдаемого состояния;
\item модуль обучения и оркестрации эксперимента, запускающий подготовку данных, обучение, тестирование и сохранение артефактов;
\item модуль тестирования и backtesting, преобразующий сигналы и действия в траекторию стоимости портфеля;
\item модуль расчета метрик и визуализации результатов, формирующий таблицы, графики и итоговые сравнительные отчеты.
\end{enumerate}

На рисунке~\ref{fig:architecture_experimental_stand} показана логика взаимодействия модулей. Важной особенностью стенда является то, что финансовые метрики выносятся в общий слой оценки. Это позволяет сравнивать модели, исходно обучавшиеся по разным критериям. Для ML-моделей такой критерий связан с качеством классификации, а для RL-моделей — с накопленной наградой. Тем не менее итоговое сравнение выполняется в одном пространстве показателей доходности и риска.

Поток данных в разработанном решении имеет следующий вид: исторические котировки поступают в модуль предобработки, затем на их основе формируется пространство признаков, после чего эти признаки либо подаются в baseline- и ML-модели, либо используются для построения состояния RL-среды. Далее агент или модель формирует торговое действие, на основе которого обновляется портфель, рассчитывается награда и строится стратегия на тестовом интервале. После завершения тестирования рассчитываются финансовые метрики и создаются графические артефакты, позволяющие интерпретировать результат.

\begin{figure}[H]
\centering
\includegraphics[width=0.92\textwidth]{architecture_experimental_stand.png}
\caption{Архитектура программного решения для прогнозирования и оптимизации торговых решений на основе RL}
\label{fig:architecture_experimental_stand}
\end{figure}

\subsection{Источники данных и структура OHLCV-данных}

В качестве источника рыночных данных в работе выбран сервис Yahoo Finance, доступ к которому осуществляется через библиотеку \textit{yfinance}. Такой выбор обусловлен несколькими причинами. Во-первых, источник обеспечивает достаточно удобный программный доступ к дневным котировкам широкого круга инструментов. Во-вторых, формат данных легко интегрируется в Python-пайплайн без промежуточных ручных преобразований. В-третьих, полученные данные пригодны для воспроизводимых учебно-исследовательских экспериментов [24].

Базовой единицей наблюдения является строка формата OHLCV, содержащая дату, цену открытия, максимум, минимум, цену закрытия, скорректированную цену закрытия и объем. Для торговых задач цена закрытия используется как основа расчета доходностей и большинства технических индикаторов. Скорректированная цена закрытия позволяет учитывать дивиденды и сплиты в тех сценариях, где это критично для корректного сравнения исторических данных [5, 24].

Использование OHLCV-данных удобно и с точки зрения расширяемости исследования. Даже если в базовом эксперименте часть признаков вычисляется только по цене закрытия, наличие полного OHLCV-набора позволяет в дальнейшем добавлять диапазонные характеристики, свечные паттерны, индикаторы объема и более сложные рыночные признаки.

\subsection{Предобработка данных}

Предобработка данных включает несколько обязательных этапов. Сначала выполняется проверка полноты схемы и корректности дат. Затем данные сортируются в хронологическом порядке и очищаются от строк с пропусками в критичных колонках. После этого производится приведение числовых полей к единому формату и отбрасывание промежуточных значений, способных нарушить воспроизводимость результатов.

Следующим важным шагом является разбиение данных на обучающую, валидационную и тестовую части без перемешивания. Для финансовых временных рядов это принципиально важно, поскольку случайное перемешивание приводит к утечке информации из будущего в прошлое. В настоящей работе используется временное разбиение, при котором сначала выделяется обучающий интервал, затем валидационный, а затем тестовый период. Такое разбиение ближе к реальному сценарию, где модель всегда обучается на исторических данных и применяется к более позднему участку ряда [2, 3, 7].

Отдельно необходимо отметить роль транзакционных издержек. Даже если в исходном ряду цены отражают рыночное движение корректно, стратегия, активно меняющая позицию, может терять значительную часть результата на комиссии. Поэтому уже на этапе проектирования стенда в торговую логику baseline-стратегий и RL-среды закладывается параметр \textit{transaction\_cost}, позволяющий избежать чрезмерно оптимистичных оценок [6, 7].

\subsection{Формирование признаков: доходности, скользящие средние, RSI, MACD, Bollinger Bands, волатильность}

Для преобразования исходных OHLCV-данных в признаки, пригодные для обучения моделей, используется набор технических и статистических характеристик, отражающих как локальную доходность, так и структуру тренда и риска. Базовыми признаками являются простая доходность и логарифмическая доходность [3, 5, 7]:
\begin{equation}
r_t = \frac{C_t}{C_{t-1}} - 1,
\end{equation}
\begin{equation}
\ell_t = \ln\left(\frac{C_t}{C_{t-1}}\right),
\end{equation}
где $C_t$ — цена закрытия в момент времени $t$.

Для описания трендовой компоненты используются скользящие средние с окнами 5, 10 и 20 периодов:
\begin{equation}
\text{SMA}_n(t) = \frac{1}{n} \sum_{i=0}^{n-1} C_{t-i}.
\end{equation}
Такой набор позволяет одновременно учитывать краткосрочную и среднесрочную динамику. Для оценки риска добавляются скользящие стандартные отклонения доходности на окнах 10 и 20 периодов, характеризующие локальную волатильность.

Индикатор RSI используется для оценки локальной перекупленности и перепроданности. MACD и его сигнальная линия позволяют отслеживать расхождение между быстрыми и медленными экспоненциальными средними, а полосы Боллинджера помогают оценить отклонение текущей цены от локального среднего уровня с учетом волатильности. Выбор именно этих признаков обусловлен их широкой представленностью в практических торговых исследованиях и достаточной интерпретируемостью [5, 7].

Все признаки рассчитываются только на исторической информации, доступной на текущем шаге. После генерации признаки проходят очистку от значений \textit{NaN}, возникающих из-за оконных вычислений. Благодаря этому обучающие алгоритмы получают согласованную матрицу признаков без утечки информации из будущего.

\subsection{Формирование целевых переменных для ML-моделей}

Для классических моделей машинного обучения в работе используются две целевые переменные. Первая — \textit{target\_direction}, представляющая собой бинарный индикатор направления изменения цены на следующем шаге [7]:
\begin{equation}
y_t =
\begin{cases}
1, & \text{если } C_{t+1} > C_t,\\
0, & \text{иначе.}
\end{cases}
\end{equation}
Эта постановка удобна для задач классификации и позволяет оценивать качество моделей по стандартным ML-метрикам.

Вторая целевая переменная — \textit{target\_return}, то есть доходность следующего периода:
\begin{equation}
\rho_t = \frac{C_{t+1}}{C_t} - 1.
\end{equation}
Хотя в основном экспериментальном контуре работа сосредоточена на классификации направления движения, использование доходности в качестве дополнительной переменной важно для интерпретации результата. Во-первых, она позволяет перевести прогноз модели в торговый эффект. Во-вторых, ее наличие упрощает последующую конвертацию предсказаний в стратегию и расчет совокупной доходности.

\subsection{Baseline-стратегии: Buy \& Hold и Moving Average Crossover}

В качестве простейших аналогов в стенд включены две baseline-стратегии. Первая стратегия — \textit{Buy \& Hold} — предполагает покупку актива в начале тестового интервала и удержание позиции до его завершения. Эта стратегия особенно важна в задачах анализа фондового рынка, поскольку на долгосрочных восходящих участках она может обеспечивать высокий абсолютный результат при минимальной сложности реализации. В то же время \textit{Buy \& Hold} почти не управляет риском и не адаптируется к смене рынка [5, 6].

Вторая стратегия — \textit{Moving Average Crossover} — использует пересечение короткой и длинной скользящих средних. Если короткая средняя находится выше длинной, стратегия считается находящейся в длинной позиции; иначе позиция закрывается. Такой подход отражает распространенную практику использования трендовых индикаторов и служит промежуточным эталоном между пассивной стратегией и моделями, обучаемыми на данных [5, 6].

Принципиально важно, что обе стратегии оцениваются по тем же финансовым показателям, что и ML- и RL-подходы. Это позволяет интерпретировать выигрыш сложной модели не в абстрактных терминах, а через конкретное улучшение доходности, просадки или коэффициента Шарпа относительно понятных и прозрачных альтернатив.

\subsection{Постановка задачи обучения с подкреплением}

Для RL-части исследования задача формулируется как задача последовательного принятия решений в среде, описываемой процессом принятия решений Маркова. На вход среды поступает DataFrame с историческими OHLCV-данными и рассчитанными признаками, а также список колонок, формирующих наблюдаемое состояние. На каждом шаге агент наблюдает состояние среды, выбирает одно из допустимых действий и получает награду, зависящую от изменения стоимости портфеля и комиссий за сделки. Таким образом, целью оптимизации становится не качество одношагового прогноза, а максимизация накопленного полезного результата на всем тестовом интервале [1].

В настоящей работе рассматривается long-only-постановка с дискретными действиями \textit{hold}, \textit{buy} и \textit{sell}. Такое ограничение упрощает интерпретацию и делает пространство решений совместимым с DQN-алгоритмом. Одновременно эта же среда может использоваться с PPO, что обеспечивает сопоставимость value-based и policy-based подходов [8, 9]. В реализованной среде агент не открывает короткие позиции, а операция \textit{buy} использует весь доступный денежный остаток, тогда как операция \textit{sell} полностью закрывает текущую длинную позицию.

Выбор именно такой постановки продиктован прикладной логикой разработанного прототипа. Во многих прикладных задачах достаточно ответить на вопрос, стоит ли в данный момент сохранять позицию, открыть ее или выйти в кэш. Даже такой ограниченный набор действий уже позволяет оценить, в какой степени RL-агент способен использовать признаки рынка для адаптивной торговли и дает возможность сравнить его с более простыми реализованными альтернативами.

\subsection{Описание состояния, действий и функции вознаграждения RL-агента}

Состояние RL-агента формируется как окно последних \textit{window\_size} наблюдений по выбранным признакам, дополненное информацией о текущей позиции и доле капитала, находящейся в наличных средствах. В реализованной версии среды \texttt{window\_size} по умолчанию равно 30, а в набор признаков входят \textit{simple\_return}, \textit{log\_return}, \textit{rolling\_mean\_5}, \textit{rolling\_mean\_10}, \textit{rolling\_mean\_20}, \textit{rolling\_std\_10}, \textit{rolling\_std\_20}, \textit{rsi\_14}, \textit{macd}, \textit{macd\_signal}, \textit{bollinger\_upper} и \textit{bollinger\_lower}. Такой вектор состояния совмещает сведения о локальной рыночной динамике и внутреннем состоянии портфеля. Агент не получает доступа к будущим значениям и действует только на основе исторической информации [1, 8, 9].

Как показано на рисунке~\ref{fig:rl_environment_scheme}, действие \textit{buy} переводит свободный капитал в позицию по активу, действие \textit{sell} закрывает позицию и возвращает капитал в денежную форму, а действие \textit{hold} сохраняет текущую структуру портфеля. Формально пространство действий можно описать следующим образом: 0 --- удержание текущего состояния портфеля, 1 --- покупка актива на весь доступный объем денежных средств, 2 --- продажа всей текущей позиции. Вознаграждение определяется через изменение \textit{portfolio\_value} между соседними шагами с учетом уплаченной комиссии. Такой выбор reward function делает обучение более прикладным, поскольку агент оптимизирует непосредственно то, что имеет финансовый смысл.

\begin{figure}[H]
\centering
\includegraphics[width=0.9\textwidth]{rl_environment_scheme.png}
\caption{Схема взаимодействия RL-агента со средой}
\label{fig:rl_environment_scheme}
\end{figure}

При этом следует учитывать, что даже разумно выбранная функция вознаграждения не гарантирует универсально качественного поведения. Если агент получает награду только за локальный прирост капитала, он может оказаться склонен к чрезмерной активности. Если, напротив, reward чрезмерно штрафует риск, агент может выбирать слишком осторожную политику. Поэтому в дипломной работе RL рассматривается как чувствительный к постановке задачи метод, а не как универсальная замена прогнозным моделям. Ограничениями разработанного решения на данном этапе являются long-only-логика, торговля одним активом в каждый момент времени, фиксированная модель комиссии и отсутствие моделирования проскальзывания.

\subsection{Метрики оценки качества моделей и стратегий}

Поскольку в стенде сравниваются модели разной природы, для оценки используются две группы метрик: классификационные и финансовые. Классификационные метрики --- Accuracy, Precision, Recall, F1-score и ROC-AUC --- применяются для ML-моделей, прогнозирующих направление движения. Они позволяют понять, насколько корректно модель различает положительный и отрицательный классы. Однако сами по себе эти показатели не отражают экономический результат использования модели [15].

Финансовые метрики, напротив, оценивают уже поведение торговой стратегии. Совокупная доходность определяется как [27, 28]:
\begin{equation}
R_{\text{total}} = \prod_{t=1}^{T}(1 + r_t^{(s)}) - 1,
\end{equation}
где $r_t^{(s)}$ — доходность стратегии в период $t$. Годовая волатильность, коэффициент Шарпа и максимальная просадка позволяют оценить риск, устойчивость и глубину неблагоприятных периодов:
\begin{equation}
\text{Sharpe} = \frac{\overline{r^{(s)} - r_f}}{\sigma(r^{(s)})}\sqrt{K},
\end{equation}
где $r_f$ — безрисковая ставка, $\sigma(r^{(s)})$ — стандартное отклонение доходности стратегии, $K$ — число периодов в году [27].

\begin{table}[H]
\caption{Метрики, применяемые в исследовании}
\label{tab:metrics_overview}
\begin{center}
\small
\begin{longtable}{|p{3.4cm}|p{3.2cm}|p{7.2cm}|}
\hline
\textbf{Метрика} & \textbf{Тип} & \textbf{Интерпретация} \\
\hline
\endfirsthead
\hline
\textbf{Метрика} & \textbf{Тип} & \textbf{Интерпретация} \\
\hline
\endhead
Accuracy & ML & Доля правильно классифицированных направлений движения цены \\
\hline
Precision & ML & Доля верных положительных сигналов среди всех положительных прогнозов \\
\hline
Recall & ML & Способность модели обнаруживать восходящие движения среди всех фактических восходящих движений \\
\hline
F1-score & ML & Гармоническое среднее Precision и Recall, полезно при неидеальном балансе классов \\
\hline
ROC-AUC & ML & Качество ранжирования положительного класса по вероятности \\
\hline
Total Return & Финансовая & Совокупный результат стратегии за весь тестовый период \\
\hline
Annualized Return & Финансовая & Доходность, приведенная к годовому масштабу \\
\hline
Volatility & Финансовая & Годовая волатильность доходности стратегии как характеристика риска \\
\hline
Sharpe Ratio & Финансовая & Соотношение избыточной доходности и риска; показывает эффективность на единицу риска \\
\hline
Max Drawdown & Финансовая & Максимальная историческая просадка капитала от локального максимума \\
\hline
Win Rate & Финансовая & Доля прибыльных торговых периодов или сделок \\
\hline
Number of Trades & Финансовая & Интенсивность торговли и косвенный индикатор транзакционных издержек \\
\hline
\end{longtable}
\end{center}

\end{table}

Использование единого набора метрик, представленного в таблице~\ref{tab:metrics_overview}, является ключевым условием корректного сравнения. В реализованном программном обеспечении для торговых стратегий рассчитываются \textit{total\_return}, \textit{annualized\_return}, \textit{annualized\_volatility}, \textit{sharpe\_ratio}, \textit{max\_drawdown}, \textit{win\_rate}, \textit{number\_of\_trades} и \textit{calmar\_ratio}, а для ML-моделей --- Accuracy, Precision, Recall, F1-score и ROC-AUC. Отдельная регрессионная задача прогноза цены в основной программный контур не включалась, поэтому метрики MSE и MAE в базовом эксперименте не рассчитывались. Именно на этом уровне проявляется основная методологическая идея работы: модель должна оцениваться не только по тому, насколько хорошо она прогнозирует цену, но и по тому, насколько полезны ее решения с точки зрения итоговой динамики капитала.

\subsection{Выводы по главе}

Во второй главе была сформирована методическая основа исследования. Определена архитектура экспериментального стенда, обоснован выбор OHLCV-данных как исходного представления рынка, описаны этапы предобработки, набор признаков и способ формирования целевых переменных для ML-моделей. Также была задана RL-постановка с дискретными действиями и функцией вознаграждения, напрямую связанной с изменением стоимости портфеля.

Ключевым результатом главы является переход от общей идеи «применить RL к котировкам» к воспроизводимой схеме эксперимента, в которой baseline-стратегии, ML-модели и RL-агенты сравниваются в едином контуре. Это создает основу для следующей главы, где описывается программная реализация построенного стенда и структура соответствующего проекта на Python.
